(1) 密度 \(\rho\)、表面张力 \(\sigma\) 的不可压无黏液体占据平底水槽 \(-h<z<0\)，重力 \(g\)，水面小波。设势流且槽底不可穿透，构造满足Laplace方程的线性波动势。

(2) 由水面线性运动、动力边界条件与表面张力求色散关系，给出深水和浅水两极限。

(3) 在深水极限求相速度极小值及相应波数，证明在该点群速度等于相速度。

(4) 讨论零表面张力、浅水长波群速度以及高波数理论对黏性和波陡限制的敏感性。
